\documentclass[12pt,a4paper]{article}

\usepackage[margin=0.9in]{geometry}
\usepackage{amsmath,amssymb}
\usepackage{booktabs}
\usepackage{array}
\usepackage{longtable}
\usepackage{enumitem}
\usepackage{hyperref}
\usepackage{xcolor}
\usepackage{fancyhdr}

% Neutral palette for formal academic presentation
\definecolor{primary}{RGB}{0,0,0}
\definecolor{secondary}{RGB}{60,60,60}
\definecolor{accent}{RGB}{85,85,85}
\definecolor{highlight}{RGB}{110,110,110}
\definecolor{danger}{RGB}{140,140,140}
\definecolor{light}{RGB}{245,245,245}

\hypersetup{
    hidelinks,
    pdftitle={Artificial Intelligence: Learning Algorithms, Regression, and Symbolic AI},
    pdfauthor={Department of Computing}
}

\pagestyle{fancy}
\fancyhf{}
\lhead{Artificial Intelligence Notes}
\rhead{Learning and Symbolic AI}
\rfoot{Page \thepage}

\setlength{\parskip}{0.5em}
\setlength{\parindent}{0pt}
\setlength{\headheight}{15pt}

\begin{document}

\begin{titlepage}
\centering
\vspace*{0.7cm}
{\Large\color{primary}\textbf{Academic Report}}\\[0.3cm]
\rule{\textwidth}{1.6pt}\\[0.5cm]
{\Huge\bfseries\color{primary}Artificial Intelligence}\\[0.2cm]
{\Large Learning Algorithms, Linear Regression, and Symbolic AI}\\[0.4cm]
\rule{\textwidth}{1.6pt}\\[0.9cm]

\begin{tabular}{rl}
\textbf{Course Focus:} & AI Fundamentals and Machine Learning \\
\textbf{Coverage:} & Supervised, Unsupervised, Semi-Supervised Learning \\
\textbf{} & Gradient Descent, Regression, Symbolic AI, Turing Test \\
\textbf{Prepared For:} & Formal Academic Review with Worked Examples
\end{tabular}

\vfill
{\large Department of Computing}\\
{\large \today}
\end{titlepage}

\tableofcontents
\newpage

\section{\textcolor{primary}{Introduction to AI and Contemporary Applications}}

Artificial Intelligence (AI) is the field of creating systems that can perform tasks that usually require human intelligence, such as perception, decision-making, language understanding, and learning from data.

\textbf{Contemporary applications of AI include:}
\begin{itemize}[noitemsep]
    \item Healthcare diagnostics and medical image analysis
    \item Recommendation systems (e.g., streaming and e-commerce)
    \item Fraud detection in banking
    \item Autonomous driving and driver assistance
    \item Smart assistants and conversational systems
    \item Predictive maintenance in industry
\end{itemize}

\textbf{AI in real life:} AI is often embedded in everyday tools, from spam filtering to route optimization in maps.

\section{\textcolor{primary}{History of AI, Controversies, and Strong vs Weak AI}}

\subsection*{Brief history}
\begin{itemize}[noitemsep]
    \item \textbf{1950s--1960s:} Early optimism, symbolic reasoning, and logic-based systems.
    \item \textbf{1970s--1980s:} AI winters due to limited compute and unmet expectations.
    \item \textbf{1990s--2010s:} Statistical learning and data-driven AI rise.
    \item \textbf{2010s--present:} Deep learning breakthroughs and large-scale AI deployment.
\end{itemize}

\subsection*{AI controversies}
Key concerns include bias, explainability, privacy, surveillance, misuse, and job displacement.

\subsection*{Strong AI vs Weak AI}
\begin{itemize}[noitemsep]
    \item \textbf{Weak AI (Narrow AI):} Built for specific tasks (current practical systems).
    \item \textbf{Strong AI (General AI):} Hypothetical human-level general intelligence across domains.
\end{itemize}

\section{\textcolor{primary}{Categories of Learning Algorithms}}

\subsection*{1) Supervised Learning}
Model learns from labeled data $(x, y)$.

\textbf{Typical tasks:}
\begin{itemize}[noitemsep]
    \item \textbf{Regression:} Predicting continuous values (house price, sales forecast, temperature).
    \item \textbf{Classification:} Predicting categorical labels (spam/non-spam, disease/no disease).
    \item \textbf{Ranking:} Ordering items by relevance (search engines, recommendation feeds).
\end{itemize}

\textbf{Common algorithms:}
\begin{itemize}[noitemsep]
    \item Linear Regression, Logistic Regression
    \item Decision Trees, Random Forests, Gradient Boosting
    \item Support Vector Machines, Neural Networks
\end{itemize}

\subsection*{2) Unsupervised Learning}
Model learns patterns from unlabeled data $x$.

\textbf{Typical tasks:}
\begin{itemize}[noitemsep]
    \item \textbf{Clustering:} Grouping similar records (customer segmentation with K-Means, hierarchical clustering).
    \item \textbf{Dimensionality reduction:} Compressing feature space while preserving structure (PCA, t-SNE, UMAP).
    \item \textbf{Association mining:} Discovering co-occurrence rules (market basket analysis).
    \item \textbf{Anomaly detection:} Identifying rare or abnormal samples.
\end{itemize}

\textbf{Why it matters:}
\begin{itemize}[noitemsep]
    \item Useful for exploratory data analysis when labels are unavailable.
    \item Helps feature engineering before supervised learning.
\end{itemize}

\subsection*{3) Semi-Supervised Learning}
Uses a small labeled set and a large unlabeled set.

\textbf{Use case:} When labeling is expensive (medical images, legal documents, satellite data, speech transcripts).

\textbf{Common strategies:}
\begin{itemize}[noitemsep]
    \item \textbf{Self-training / pseudo-labeling:} Train on labeled data, label high-confidence unlabeled data, retrain.
    \item \textbf{Consistency regularization:} Model should produce stable predictions under small input perturbations.
    \item \textbf{Graph-based methods:} Nearby points in feature space are encouraged to share labels.
\end{itemize}

\textbf{Example:} Train initial classifier on labeled emails, then pseudo-label high-confidence unlabeled emails and retrain.

\subsection*{4) Reinforcement Learning (contextual category)}
Although not in the previous three, reinforcement learning is another major category.
An agent interacts with an environment and learns a policy to maximize cumulative reward.
\begin{itemize}[noitemsep]
    \item Examples: game playing, robotics control, resource scheduling.
    \item Core terms: state, action, reward, policy, value function.
\end{itemize}

\section{\textcolor{primary}{Datasets, Variables, and Features}}

\begin{itemize}[noitemsep]
    \item \textbf{Dataset:} Collection of observations/records.
    \item \textbf{Variable:} Any measurable quantity in the dataset.
    \item \textbf{Feature:} Input variable used by model to predict output.
    \item \textbf{Target/Label:} Output variable to be predicted in supervised learning.
\end{itemize}

\textbf{Simple example:} Predict salary from years of experience.
\begin{itemize}[noitemsep]
    \item Feature: Years of experience
    \item Target: Salary
\end{itemize}

\section{\textcolor{primary}{Linear Regression and Regression Model}}

Linear regression models relationship between target $y$ and one or more features $x$.

\subsection*{Simple Linear Regression}
\begin{equation}
\hat{y} = w x + b
\end{equation}
where $w$ is weight (slope) and $b$ is bias (intercept).

\subsection*{Multiple Linear Regression}
\begin{equation}
\hat{y} = w_1x_1 + w_2x_2 + \cdots + w_nx_n + b
\end{equation}

\textbf{Model parameter estimation:} Find $w$ and $b$ values that minimize prediction error.

\section{\textcolor{primary}{Loss Function, Cost Function, and Optimization Terms}}

\textbf{Loss function:} Error for a single training example.
\begin{equation}
L^{(i)} = (\hat{y}^{(i)} - y^{(i)})^2
\end{equation}

\textbf{Cost function:} Average loss over all $m$ examples.
\begin{equation}
J(w,b) = \frac{1}{2m}\sum_{i=1}^{m}(\hat{y}^{(i)} - y^{(i)})^2
\end{equation}

\textbf{Objective function:} General function to minimize or maximize.

\textbf{Optimization problem:}
\begin{itemize}[noitemsep]
    \item \textbf{Unconstrained:} Minimize $J(\theta)$ without restrictions.
    \item \textbf{Constrained:} Minimize $J(\theta)$ subject to constraints (e.g., $g(\theta) \le 0$).
\end{itemize}

\section{\textcolor{primary}{Gradient Descent and Gradient Ascent}}

\subsection*{Gradient Descent (Minimization)}
Used to minimize cost function:
\begin{equation}
w \leftarrow w - \alpha \frac{\partial J}{\partial w}, \quad b \leftarrow b - \alpha \frac{\partial J}{\partial b}
\end{equation}
where $\alpha$ is learning rate.

\subsection*{Gradient Ascent (Maximization)}
Used to maximize objective function:
\begin{equation}
\theta \leftarrow \theta + \alpha \nabla f(\theta)
\end{equation}

\subsection*{Why descent works}
Gradient points in direction of steepest increase, so subtracting it moves toward lower cost.

\section{\textcolor{primary}{Simple Linear Regression Cost Using MSE}}

For simple regression, prediction is $\hat{y}^{(i)} = wx^{(i)} + b$, and MSE-based cost is:
\begin{equation}
J(w,b) = \frac{1}{2m}\sum_{i=1}^{m}(wx^{(i)} + b - y^{(i)})^2
\end{equation}

Gradients:
\begin{equation}
\frac{\partial J}{\partial w} = \frac{1}{m}\sum_{i=1}^{m}(wx^{(i)} + b - y^{(i)})x^{(i)}, \quad
\frac{\partial J}{\partial b} = \frac{1}{m}\sum_{i=1}^{m}(wx^{(i)} + b - y^{(i)})
\end{equation}

\section{\textcolor{primary}{Worked Example: Iterations of Gradient Descent}}

Given dataset points: $(1,2)$, $(2,3)$, $(3,4)$. Initialize $w=0$, $b=0$, learning rate $\alpha = 0.1$.

\subsection*{Iteration 0 (initial)}
Predictions: $[0,0,0]$\newline
Errors: $[-2,-3,-4]$
\begin{align}
\frac{\partial J}{\partial w} &= \frac{1}{3}[(-2)(1)+(-3)(2)+(-4)(3)] = -\frac{20}{3} \approx -6.6667 \\
\frac{\partial J}{\partial b} &= \frac{1}{3}[(-2)+(-3)+(-4)] = -3
\end{align}
Update:
\begin{equation}
w=0-0.1(-6.6667)=0.6667,\quad b=0-0.1(-3)=0.3
\end{equation}

\subsection*{Iteration 1}
Using $w=0.6667$, $b=0.3$:
\begin{itemize}[noitemsep]
    \item Predictions: $[0.9667,1.6334,2.3001]$
    \item Errors: $[-1.0333,-1.3666,-1.6999]$
\end{itemize}

New gradients:
\begin{align}
\frac{\partial J}{\partial w} &\approx -3.6221 \\
\frac{\partial J}{\partial b} &\approx -1.3666
\end{align}
Update:
\begin{equation}
w\approx1.0289,\quad b\approx0.4367
\end{equation}

\subsection*{Iteration summary table}
\begin{center}
\begin{tabular}{cccc}
\toprule
Iteration & $w$ & $b$ & Cost $J(w,b)$ (approx) \\
\midrule
0 & 0.0000 & 0.0000 & 4.8333 \\
1 & 0.6667 & 0.3000 & 1.0123 \\
2 & 1.0289 & 0.4367 & 0.3168 \\
3 & 1.2255 & 0.4933 & 0.1770 \\
\bottomrule
\end{tabular}
\end{center}

\textbf{Observation:} Cost decreases each iteration, demonstrating minimization.

\section{\textcolor{primary}{Minimization of Cost Function: Practical Notes}}

\begin{itemize}[noitemsep]
    \item If $\alpha$ is too small, convergence is very slow.
    \item If $\alpha$ is too large, updates may overshoot and diverge.
    \item Feature scaling often improves convergence speed.
\end{itemize}

\section{\textcolor{primary}{Explained and Unexplained Variation}}

Let mean of target be $\bar{y}$.

\textbf{Total deviation / total sum of squares (SST):}
\begin{equation}
SST = \sum_{i=1}^{m}(y_i - \bar{y})^2
\end{equation}

\textbf{Unexplained variation / sum of squared errors (SSE):}
\begin{equation}
SSE = \sum_{i=1}^{m}(y_i - \hat{y}_i)^2
\end{equation}

\textbf{Explained variation (SSR):}
\begin{equation}
SSR = \sum_{i=1}^{m}(\hat{y}_i - \bar{y})^2
\end{equation}

For linear regression with intercept: $SST = SSR + SSE$.

\subsection*{Intuition}
\begin{itemize}[noitemsep]
    \item $SST$ tells us total spread of observed values around the mean.
    \item $SSR$ is the part of spread captured by the model (signal explained by predictors).
    \item $SSE$ is leftover spread not captured by the model (noise, missing factors, model mismatch).
\end{itemize}

\subsection*{Numerical example}
Suppose observed values are $y=[3,5,7]$, predicted values are $\hat{y}=[2.5,5.5,7]$.
Then $\bar{y}=5$.
\begin{align}
SST &= (3-5)^2 + (5-5)^2 + (7-5)^2 = 8 \\
SSE &= (3-2.5)^2 + (5-5.5)^2 + (7-7)^2 = 0.5 \\
SSR &= (2.5-5)^2 + (5.5-5)^2 + (7-5)^2 = 10.5
\end{align}

In exact arithmetic for OLS with intercept, fitted values satisfy orthogonality and decomposition constraints.
For hand-rounded predictions or non-OLS estimates, direct computed $SSR$ may not perfectly satisfy $SST=SSR+SSE$ due to approximation effects.

\subsection*{Coefficient of determination ($R^2$)}
\begin{equation}
R^2 = 1 - \frac{SSE}{SST}
\end{equation}

Interpretation: proportion of target variance explained by the model.
\begin{itemize}[noitemsep]
    \item $R^2=1$: perfect fit.
    \item $R^2=0$: model is no better than predicting the mean.
    \item $R^2<0$: model performs worse than mean predictor (possible on test data or poorly specified models).
\end{itemize}

\subsection*{Correlation coefficient ($r$)}
For simple linear regression:
\begin{equation}
r = \frac{\sum (x_i-\bar{x})(y_i-\bar{y})}{\sqrt{\sum (x_i-\bar{x})^2\sum (y_i-\bar{y})^2}},\quad R^2 = r^2
\end{equation}

\textbf{Sign of $r$:} positive means both variables increase together; negative means one increases while the other decreases.

\section{\textcolor{primary}{Comprehensive Terminology Quick List}}

\begin{longtable}{>{\raggedright\arraybackslash}p{0.28\textwidth}p{0.66\textwidth}}
\toprule
\textbf{Term} & \textbf{Meaning} \\
\midrule
Weights & Parameters multiplying input features (importance of each feature). \\
Bias & Constant term shifting prediction function up or down. \\
Loss Function & Error for one example. \\
Cost Function & Average (or aggregate) loss across dataset. \\
Objective Function & Function being optimized (minimized or maximized). \\
Optimization & Process of finding best parameter values for objective. \\
Constraints & Restrictions on feasible parameter values. \\
Constrained Optimization & Optimization with explicit restrictions. \\
Unconstrained Optimization & Optimization without explicit restrictions. \\
Regression Model & Mathematical mapping from features to continuous target. \\
\bottomrule
\end{longtable}

\section{\textcolor{primary}{Symbolic AI and Physical Symbol System Hypothesis}}

\subsection*{Symbolic AI}
Symbolic AI represents knowledge using symbols and explicit rules (if-then logic), unlike pure statistical pattern matching.

\textbf{How symbolic AI represents knowledge:}
\begin{itemize}[noitemsep]
    \item \textbf{Logic formulas:} Facts and rules in propositional or predicate logic.
    \item \textbf{Production rules:} IF condition THEN action.
    \item \textbf{Semantic networks / ontologies:} Concepts and relations.
    \item \textbf{Frames/scripts:} Structured templates for stereotyped situations.
\end{itemize}

\textbf{Reasoning styles:}
\begin{itemize}[noitemsep]
    \item \textbf{Deductive reasoning:} General rule to specific conclusion.
    \item \textbf{Abductive reasoning:} Best explanation for observed evidence.
    \item \textbf{Rule chaining:} Forward chaining (data-driven) and backward chaining (goal-driven).
\end{itemize}

\textbf{Example (rule-based medical triage):}
\begin{itemize}[noitemsep]
    \item Rule 1: IF fever \texttt{\&} cough THEN possible respiratory infection.
    \item Rule 2: IF respiratory infection \texttt{\&} low oxygen THEN urgent referral.
\end{itemize}
Given facts \{fever, cough, low oxygen\}, the inference engine derives urgent referral.

\textbf{Strengths of symbolic AI:}
\begin{itemize}[noitemsep]
    \item Transparent and interpretable decision process.
    \item Easy to encode domain knowledge from experts.
    \item Works well in domains with clear logic and constraints (law, configuration, planning).
\end{itemize}

\textbf{Limitations:}
\begin{itemize}[noitemsep]
    \item Knowledge acquisition bottleneck (manual rule engineering is expensive).
    \item Fragility in noisy, ambiguous, real-world perception tasks.
    \item Limited scalability compared to modern data-driven learning for high-dimensional inputs.
\end{itemize}

\textbf{Neuro-symbolic direction:} Hybrid systems combine neural perception with symbolic reasoning for better robustness and explainability.

\subsection*{Physical Symbol System Hypothesis}
Proposed by Newell and Simon: a physical symbol system has the necessary and sufficient means for general intelligent action.

\textbf{Core idea:} Intelligence can emerge from symbol manipulation through formal rules.

\textbf{Interpretation in AI history:}
\begin{itemize}[noitemsep]
    \item Motivated early expert systems and theorem provers.
    \item Emphasized cognition as computation over symbolic structures.
    \item Sparked ongoing debate with connectionist approaches (neural networks).
\end{itemize}

\textbf{Modern view:} Symbol manipulation is powerful for reasoning and structure, while sub-symbolic learning is powerful for perception; practical AI often benefits from combining both.

\section{\textcolor{primary}{Classical Chatbots and the Turing Test}}

\subsection*{Classical chatbots}
\begin{itemize}[noitemsep]
    \item \textbf{ELIZA:} Pattern-matching therapist-style responses.
    \item \textbf{ALICE:} Rule-based chatbot using AIML.
    \item \textbf{Kuki (Mitsuku):} Competition-winning conversational bot with scripted/rule components.
\end{itemize}

\subsection*{Turing Test}
Proposed by Alan Turing (1950): if a human evaluator cannot reliably distinguish machine from human in text conversation, machine is considered intelligent under this test.

\subsection*{Common objections to Turing Test}
\begin{itemize}[noitemsep]
    \item Passing conversation test does not prove understanding or consciousness.
    \item Systems may mimic language without real semantic grounding.
    \item It measures imitation, not necessarily general intelligence.
\end{itemize}

\section{\textcolor{primary}{Conclusion}}

This document reviewed key AI foundations, major learning paradigms, and detailed linear regression parameter estimation with gradient descent. It also covered optimization vocabulary, variation-based evaluation metrics ($SSE$, $SST$, $R^2$, correlation), and foundational symbolic AI concepts including classical chatbots and the Turing test.

\textbf{Revision note:} Combine precise definitions with one worked numerical example (e.g., gradient descent iterations) and one conceptual critique (e.g., limits of the Turing test).

\end{document}
